\chapter{The Landscape}
\label{ch:landscape}

\begin{quote}\itshape
Most of what you will find when you search for this is describing a library that
no longer exists.
\end{quote}

\noindent
This chapter is about knowing which decade you are in. LangChain 1.0 landed on
22~October~2025 and did not merely add features --- it removed the three
abstractions that every tutorial written before it is built on, and inverted the
relationship between the two libraries whose names are on the cover of this book.
Getting that inversion the right way round is the single highest-value thing in
the chapter, because everything downstream follows from it.

\mermaidfig{lc-lg-layering}{%
  The inversion, drawn. \api{create\_agent} is not an alternative to LangGraph;
  it returns a compiled LangGraph graph, which is why an agent gets
  checkpointing, streaming and interrupts without asking for them.}{%
  Layer cake: application, convenience layer, runtime, contracts.}

\chapterstub{
\textbf{Argument.} \pkg{langchain} and \pkg{langgraph} are not alternatives.
\pkg{langgraph} is the execution runtime; \pkg{langchain} 1.x is a thin,
opinionated layer over it, and \api{create\_agent} returns a compiled LangGraph
graph. Everything an agent gets for free --- checkpointing, streaming,
interrupts --- it gets because of that fact.

\textbf{Sections.} (1.1) Timeline, from October 2022 to the present, and what
each date changed. (1.2) What 1.0 removed and why: \api{AgentExecutor},
\api{initialize\_agent}, \api{LLMChain}, \code{langgraph.prebuilt.create\_react\_agent}.
(1.3) The package map --- what each distribution owns and what moved to
\pkg{langchain-classic}. (1.4) The inversion, demonstrated rather than asserted:
build an agent with \api{create\_agent}, then print its graph. (1.5) How to date
a piece of LangChain material in ten seconds. (1.6) Where the framework is the
wrong tool at all --- forward reference to Chapter~\ref{ch:ecosystem}.
(1.7) The .NET reader's orientation: what maps, what does not, and the honest
account of what a decade of backend work is worth here.

\textbf{Listings.} Package map as a dependency listing; the before/after pair
(0.x \api{AgentExecutor} vs 1.x \api{create\_agent}) shown side by side; the
one-liner that prints the compiled graph of an agent.

\textbf{Diagrams.} \code{lc-lg-layering} --- the layer cake, showing
\api{create\_agent} sitting on the Pregel runtime. \code{lc-timeline} --- the
ecosystem timeline as a Mermaid timeline. \code{lc-package-map} --- the
distributions and their dependency edges.

\textbf{Source topics covered.} Part~III 3.1--3.2, 3.16 (partial);
Part~VIII 8.6; \S0 framing.
}
